\emph{(i)} By Premise \claimref{singularity.c1-curvature-cutoff}, as used in Proposition \claimref{singularity.stretched-singularity-count}, the locus is $2M/r^{3} = 1$, $r_{c}^{3} = 2M$. Then $r_{c} < 2M$ if and only if $r_{c}^{3} < (2M)^{3}$, that is $2M < 8M^{3}$, that is $M^{2} > \tfrac14$; at $M = \tfrac12$ the locus is the horizon, below it lies outside. On $r = 2M$ the scale is $2M/(2M)^{3} = 1/(4M^{2}) < 1$ for $M > \tfrac12$, and since $2M/r^{3}$ decreases in $r$, the whole exterior $r \ge 2M$ is below the Planck value. For Kerr with $a \ge 2(2M)^{2/3}$, Theorem \claimref{synthesis.kerr-ring-count-misses-area-law}(1) places the locus in $|r| \le \rho_{c} < r_{-}$ and shows $|z| \ge r_{-} > \rho_{c}$ for $r_{-} \le r \le r_{+}$; on and outside $r = r_{+}$, $|z| \ge r_{+}$ and the scale $2M|z|^{-3} \le 2M/r_{+}^{3} < 1$. Premises (C1) and (C2) are stated for the interior of a hole (Schwarzschild, and Kerr through Theorem \claimref{synthesis.kerr-ring-count-misses-area-law}) and place modes only on its cutoff locus; Proposition \claimref{singularity.stretched-singularity-count} counts $\nu$ modes per Planck cell of that locus. So they assign no modes to any horizon of a hole, nor to any local Rindler horizon, which is a small neighbourhood of a point off every cutoff locus. The check verifies the threshold and the horizon value at rational masses, and the Kerr inequalities $\rho_{c} < r_{-} < r_{+}$, $2M/r_{+}^{3} < 1$ at $(M, a) = (500, 300), (500, 400), (500, 480)$, where $2M = 10^{3}$ gives $\rho_{c} = 10$ and $\sqrt{M^{2} - a^{2}} = 400, 300, 140$ exactly.

\emph{(ii)} Let $\xi$ be a Killing field normal to a null hypersurface $\mathcal{N}$, and $q_{ab}$ the induced metric on cross-sections. The null extrinsic curvature of the generators is $B_{ab} = q_{a}{}^{c}q_{b}{}^{d}\nabla_{d}\xi_{c}$; its symmetric part vanishes by Killing's equation $\nabla_{(c}\xi_{d)} = 0$, and its antisymmetric part (the twist) vanishes because $\xi$ is hypersurface orthogonal on $\mathcal{N}$. The expansion and shear are the trace and the trace-free symmetric part of $B_{ab}$, so both vanish at every point of $\mathcal{N}$. Rescaling $\xi$ to an affinely parametrised $k = e^{f}\xi$ multiplies $B_{ab}$ by $e^{f}$ and adds terms proportional to $\xi_{a}$, which the projections kill, so the statement does not depend on the parametrisation (standard; e.g.\ Wald's text, cited from memory). Hence at every $p$ on a Killing horizon the cross-section through $p$ is a spacelike two-surface element whose null normal $\ell$ (the generator) has vanishing expansion and shear at $p$, which is the defining property of $\mathcal{H}_{p,\ell}$ in Premise \claimref{synthesis.c-loc-local-horizon-premise} and in Fact \claimref{synthesis.imported-jacobson-clausius-1995}; the horizon agrees with $\mathcal{H}_{p,\ell}$ at leading order in the patch size. The check computes the expansion $\theta = \tfrac12 f\,(r^{2})'/r^{2} = f/r$ of the outgoing radial null field $\partial_{v} + \tfrac12 f\,\partial_{r}$ in ingoing Eddington-Finkelstein coordinates, $f = 1 - 2M/r$, as a rational function, and confirms $\theta = 0$ at $r = 2M$ and $\theta \ne 0$ at $r = 3M$; the shear vanishes there by spherical symmetry. Under (H) each patch has entropy $\eta$ per unit area, so a cross-section of area $A$ has $\eta A$; at $\eta = \tfrac14$ this is $A/4$, which for Schwarzschild is $\SBH$ and for Kerr is $A_{+}/4$, equal to the entropy of Fact \claimref{synthesis.imported-kerr-surface-gravity-entropy} at every $a < M$.

\emph{(iii)} By Theorem \claimref{scrambling.epoch-modular-period-two-unknowns-become}, as used in Theorem \claimref{synthesis.clausius-route-newton-constant-is-residue}(ii), the epoch count of the cutoff modes is $4\pi\nu A$ modes, entropy $4\pi\nu\langle s\rangle A$. By (ii) the horizon modes carry $\eta A = 4\pi\nu\langle s\rangle A$ with the same constants (that is what ``with the constants of the Killing case'' in (H) says). For a product state of two sets of modes the entropy is the sum, so the hole's boundary modes carry $8\pi\nu\langle s\rangle A$. Theorem \claimref{scrambling.area-law-test} reads the count of the hole's boundary modes as $\SBH$, and Theorem \claimref{synthesis.quarter-from-temperature} gives the hole's thermodynamic entropy $\SBH = \int dM/T = 4\pi M^{2} = A/4$. Under (S$_{\mathrm{tot}}$), this is the total over all the hole's boundary modes; so, $8\pi\nu\langle s\rangle = \tfrac14$, so $\pi\nu\langle s\rangle = \tfrac1{32}$ and $\eta = 4\pi\nu\langle s\rangle = \tfrac18$. Fact \claimref{synthesis.imported-jacobson-clausius-1995} gives $G_{\mathrm{cl}} = 1/(4\eta) = 2$, in units of $G$. With the registered $\pi\nu\langle s\rangle = \tfrac1{16}$ kept, the sum is $8\pi\nu\langle s\rangle A = A/2$. Without (S$_{\mathrm{tot}}$) the registered value stands, $G_{\mathrm{cl}} = G$, and the conclusion is only that the boundary modes carry twice the thermodynamic entropy. With $k$ independent sets, $k$ counting all of them, the cutoff set included, $4k\pi\nu\langle s\rangle = \tfrac14$ and $G_{\mathrm{cl}} = k$. The check carries $x = \pi\nu\langle s\rangle$ as a rational number through these steps for $k = 1, 2, 3$.

\emph{(iv)} Under identification the hole's single set of modes has the entropy (H) assigns, $A_{+}/4$ by (ii), and the entropy that the literal count assigns, which in the 1-nat convention is the number $N$ of Theorem \claimref{synthesis.kerr-ring-count-misses-area-law}(4): $N/(A_{+}/4) = (\pi/2)\rho_{c}^{2}(1 - \eta_{K})/\sqrt{M^{2} - a^{2}}$ with $0 \le \eta_{K} \le \varepsilon^{2}/4 + 5\delta^{2}/32$, $\varepsilon = \rho_{c}^{2}/a^{2}$, $\delta = 2M\rho_{c}/a^{2}$ (there written $\eta$; it is not the area density). The check recomputes the ratio from the Kerr data: $\kappa_{+} = (r_{+} - r_{-})/(2(r_{+}^{2} + a^{2}))$, confirmed equal to $\sqrt{M^{2} - a^{2}}/(2Mr_{+})$, $A_{+} = 4\pi(r_{+}^{2} + a^{2}) = 8\pi Mr_{+}$, $N = \nu V'\cdot 2\pi/\kappa_{+}$ with $\nu = 1/16\pi$ and $V' = 4\pi^{2}\rho_{c}^{2}(1 - \eta_{K})$, giving $N/(A_{+}/4) = \pi c(1 - \eta_{K})$ with $c = \rho_{c}^{2}/(2\sqrt{M^{2} - a^{2}})$. With $3 < \pi < \tfrac{22}{7}$: at $(500, 300)$, $c = \tfrac18$ and the ratio is at most $\tfrac{22}{7}\cdot\tfrac18 = \tfrac{11}{28}$; at $(500, 480)$, $c = \tfrac{5}{14}$, $\delta = \tfrac{10^{4}}{230400} \le \tfrac14$, and the ratio is at least $3\cdot\tfrac{5}{14}\,(1 - \varepsilon^{2}/4 - 5\delta^{2}/32) > 1$. In both cases one set would have two different entropies. Since $A_{+}/4$ is also the entropy of Fact \claimref{synthesis.imported-kerr-surface-gravity-entropy}, this is the conflict between the literal count and that Fact which Theorem \claimref{synthesis.kerr-ring-count-misses-area-law}(4) already states without (H); the identified reading inherits it. Under (S$_{\mathrm{tot}}$) the independent reading meets it too: the hole's total would be $A_{+}/4 + N \ne A_{+}/4$ since $N > 0$.

\emph{Conclusion.} As far as their entropies go, two sets of modes are a product (iii), one set (iv), or, as the residual case, correlated with some joint state. The first contradicts, under (S$_{\mathrm{tot}}$), $G_{\mathrm{cl}} = G$ (the consistency of Theorem \claimref{synthesis.clausius-route-newton-constant-is-residue}(iii)) or the area law; the second carries the literal Kerr count's existing conflict with $A_{+}/4$; the residual case is not excluded. Refutation condition (a) of (C$_{\mathrm{loc}}$) concerns the count on local horizons; the literal Kerr count is on the cutoff tube, which by (i) is no local horizon, and its conflict with $A_{+}/4$ appears in every reading, so it bears on the cutoff placement on spinning holes and not on (C$_{\mathrm{loc}}$).
